\chapter{Conclusions and Future Work}
\label{chap:conclusions}

\par This thesis designed, implemented, and evaluated a real-time intrusion detection system for IoT networks on the CICIoT2023 dataset, delivering a reproducible offline training pipeline and an interactive web demonstration. The pipeline removes the large exact-duplicate fraction before sampling, caps each class at 200{,}000 rows, and prunes the 39 public features to 25 by variance and correlation. On that input a multi-layer perceptron was trained at binary and 8-class granularities on a stratified random split, compared against a Random Forest baseline under identical metrics, and calibrated by temperature scaling. The NEURAL IDS demonstration pairs a FastAPI backend with a React dashboard that classifies uploaded captures or CSVs in under 10~ms and is deployed for remote access.

\par Four findings stand out. The binary weighted F1 (around 0.93 for the MLP) clears the 0.85 target of NFR-2, while the 8-class figure (around 0.80 for the MLP, 0.84 for the Random Forest) sits below it, the expected price of fine-grained family discrimination from transport-layer features alone. At the family level Mirai is near-perfectly detected (0.99 recall), DDoS and DoS are heavily confused with each other as volumetric floods, and the application-layer families (Web, Reconnaissance, Brute Force) are weakest because they are distinguished by payload content the header features do not capture, a ceiling set by the data, not the classifier. The Random Forest beats the MLP by 4 to 7 points across accuracy, weighted F1, and macro F1, consistent with tabular-learning results; the MLP is kept as the serving model because its inference is a single portable call and the gap stays within target.

\par The main limitations are the header-and-metadata feature set, which caps application-layer detection; the controlled 105-device testbed, which omits multi-stage campaigns and background noise; and the random split, which leaks near-duplicate flows across the train/test boundary and so reports an upper bound on deployment performance. Global temperature scaling also leaves residual per-class miscalibration (8-class ECE 0.031 after scaling).

\par These point to clear next steps: evaluate on leakage-aware (group-wise and temporal) splits to quantify how much performance the random split inflates; replace the single global temperature with per-class scaling and per-class thresholds to tune the detection/false-alarm trade-off per family; extend the features with TLS-visible metadata to reach the encrypted and web-based attacks the transport-layer view currently misses; and deploy the serving component on a real IoT gateway with periodic retraining on misclassified flows to test the pipeline beyond the laboratory.
